% arara: pdflatex
% arara: clean: { extensions: [ aux, log ] }

\documentclass{article}
%\usepackage{paralist}
\usepackage{enumitem}
\setlist{nosep}
\setitemize[1]{label=---,itemindent=2cm}
\setenumerate[1]{label=\textcircled{\scriptsize\Alph*/},font=\sffamily}
\usepackage{layouts}

\begin{document}
\section*{Lists}
\begin{enumerate}[start=4]
    \item State the paper size by an option to the document class
    \item[6/2] Determine the margin dimensions using one of these packages:
    \begin{itemize}
        \item geometry
        \item typearea
    \end{itemize}
    \item Customize header and footer by one of these packages:
    \begin{itemize}
        \item fancyhdr
        \item scrpage-scrlayer
    \end{itemize}
\end{enumerate}
\noindent\textbf{Tweaking the line spacing:}
\begin{enumerate}[resume*]
    \item Adjust the line spacing for the whole document
    \begin{itemize}
        \item by using the setspace package and one of its options:
        \begin{itemize}[label=*]
            \item singlespacing
            \item onehalfspacing
            \item double spacing
        \end{itemize}
        \item or by the command
            \verb|\linespread{factor}|
    \end{itemize}
\end{enumerate}
\begin{description}
    \item[paralist] provides compact lists and list 
    versions that can be used within paragraphs,
    helps to customize labels and layout.
    \item[enumitem] gives control over labels and lengths in all kind of lists.
    \item[mdwlist] is useful to customize description lists,
    it even allows multi-line labels.
    It features compact lists and the capability to suspend and resume.
    \item[desclist] offers more flexibility in definition list.
    \item[multenum] produces vertical enumeration in multiple columns.
\end{description}
\newpage
\listdiagram
\end{document}
